\section{Med-PaLM 的多维人类评价}

模型链条搭好后，课程回到最重要的证据问题：Med-PaLM 是否只是在选择题上更强，还是在真实长答案的科学性、风险与可用性上也缩小了与医生的差距？接下来的九张结果页不能只看颜色高低，而要逐轴阅读。

\subsection{科学共识与推理证据}

Scientific consensus 页比较 Flan-PaLM、Med-PaLM 和 clinician answers。总体趋势是 Med-PaLM 明显改善并接近专家，但仍有差距。下一页把 comprehension、knowledge retrieval 与 reasoning 拆开；课堂问答特别指出，正确证据与错误证据不是 $1-x$ 的互补关系，一段回答可以同时包含正确事实与错误推论。

\lecturefigure{slide-35-clinical-consensus.jpg}{与科学和临床共识的一致性}{官方录像恢复幻灯片 V035；Stanford CS25 Lecture 17}

\lecturefigure{slide-36-comprehension-retrieval-reasoning.jpg}{理解、医学知识检索与推理证据}{官方录像恢复幻灯片 V036；Stanford CS25 Lecture 17}

\teachervoice{现场学生问两组 comprehension 指标是否互为补数，讲者明确回答“不”。正确理解和错误理解可以同时出现，例如回答先正确识别疾病，又在治疗建议上加入不相关推断。这是阅读复杂 human-evaluation 图时非常重要的纠错。}

\begin{knowledgebox}{读图：为什么指标可以同时为真}
长答案由许多 proposition 组成。设其中一部分正确、一部分错误，则“存在正确理解证据”和“存在错误理解证据”都可为真。用一个 accuracy 把整段文字压成对错，会丢失这种混合状态；逐轴评价的价值就在于暴露局部失败。
\end{knowledgebox}

\subsection{更长的回答也可能制造更多错误}

Incorrect or missing content 页揭示对齐的副作用：Med-PaLM 被引导生成更完整、更详细的答案，但增加的内容也扩大了出现无关或错误陈述的机会。讲者解释，某些轴上 Flan-PaLM 看似略好，并不意味着它整体更安全，可能只是因为它说得更少。评价必须同时看 completeness 与 precision。

\lecturefigure{slide-37-incorrect-missing-context.jpg}{错误内容与缺失内容的双重风险}{官方录像恢复幻灯片 V037；Stanford CS25 Lecture 17}

\begin{importantbox}{临床生成中的长度悖论}
短答案可能遗漏关键警告，长答案可能引入额外幻觉。目标不是最大化字数，而是在覆盖必要决策信息的同时控制未经支持的 proposition。工程上应把“回答长度”视作风险变量，而不是天然的质量信号。
\end{importantbox}

\subsection{伤害与偏见：错误的后果并不等价}

Harm 页同时考虑可能性与严重程度。课堂问答给出具体解释：严重伤害不只来自推荐危险治疗，也可能来自 misdiagnosis、漏掉危及生命的病情，或没有传达疾病严重性。Bias 页则检查回答是否引入不当人口统计偏见。两者都需要领域专家结合上下文判断，不能用关键词黑名单代替。

\lecturefigure{slide-38-extent-likelihood-harm.jpg}{潜在伤害的程度与发生可能性}{官方录像恢复幻灯片 V038；Stanford CS25 Lecture 17}

\lecturefigure{slide-39-potential-bias.jpg}{回答中潜在偏见的评价}{官方录像恢复幻灯片 V039；Stanford CS25 Lecture 17}

\begin{warningbox}{伤害不是“答案里有没有危险词”}
模型即使没有直接给出危险建议，也可能因漏诊、低估严重程度或错误安慰而造成延迟治疗。相反，提到癌症或死亡并不自动构成伤害；关键是内容是否与问题和证据匹配，以及它可能怎样改变用户行为。
\end{warningbox}

\teachervoice{讲者把严重 harm 的例子具体化为未识别 cancer 等致命病理，或没有适当传达严重程度。他还指出这一主题很细致，论文借助 AHRQ 风险框架；课堂图只是高层概览，不应被讲义改写成完整临床安全规范。}

\subsection{普通用户看到的是 helpfulness 与可理解性}

Layperson 结果页关注意图与帮助程度。课堂给出的近似趋势是：Flan-PaLM 回答约六成被认为 helpful，Med-PaLM 提升到约八成，而 clinician answers 接近九成。重要的不是把这些比例当作部署阈值，而是看到医学对齐能缩小沟通差距，却没有达到医生水平。

\lecturefigure{slide-40-helpfulness-utility.jpg}{普通用户评价：意图覆盖、帮助程度与效用}{官方录像恢复幻灯片 V040；Stanford CS25 Lecture 17}

\lecturefigure{slide-41-clinical-qualitative-example.jpg}{定性案例：医生、Flan-PaLM 与 Med-PaLM 的回答差异}{官方录像恢复幻灯片 V041；Stanford CS25 Lecture 17}

定性示例页应该从信息组织方式读取：医生答案往往精确、简短，但专业术语可能让患者难以理解；模型可以展开解释、补充背景，却也可能因为展开而增加错误。两者的互补点是语言转换与信息整理，而不是让模型在无人监督下替代医学判断。

\begin{knowledgebox}{读图：比较答案时先看四件事}
先看是否直接回答用户意图，再看关键医学事实是否齐全，然后检查新增细节是否有证据，最后判断是否明确给出需要专业就医的边界。流畅度只能放在这四项之后；否则最会写作的回答可能被误判为最安全的回答。
\end{knowledgebox}

定性案例也说明自动 metric 的盲区。两段答案可能共享大量词汇，却在一个关键否定词、剂量或就医建议上产生完全不同后果；反过来，医生用一句专业缩写给出正确结论，语言相似度可能很低。适合的评测应先把答案分解为临床 proposition，检查每条是否由问题和共识支持，再评价组织、语气和可行动性。若系统用于把医生笔记改写给患者，还应要求保留原始文本链接和人工确认入口，使“易懂版本”不能悄悄改写医学事实。这样，语言模型的优势被用于解释和转换，而权威来源仍可追溯。

\subsection{临床语言模型还需要吗？}

倒数第二页提出一个开放问题：当通用大模型已经很强，是否还需要临床专用语言模型（clinical LM）？课堂没有给出简单否定。小型领域模型可能在成本、隐私、本地部署或特定数据分布上有价值；大型通用模型则利用更广知识和 scaling。真正问题是在哪些数据、任务和约束下，领域预训练的边际收益超过更大通用模型。

\lecturefigure{slide-42-clinical-lm-question.jpg}{开放问题：是否仍需要专门的临床语言模型？}{官方录像恢复幻灯片 V042；Stanford CS25 Lecture 17}

\lecturefigure{slide-43-clinical-key-takeaways.jpg}{临床章节的关键结论与限制}{官方录像恢复幻灯片 V043；Stanford CS25 Lecture 17}

\teachervoice{讲者对近期临床应用给出的角色不是 autonomous physician，而是 augment physicians：把难懂的 doctor speak 转成患者更容易理解的说明，帮助医生与患者、医生与医生、医生与研究者沟通。他同时说这仍“不完美”，complementarity 才是当前机会。}

\begin{importantbox}{临床章节的证据结论}
课堂证据支持三点：大型通用模型能在医学选择题上强迁移；医学 prompt tuning 能改善长答案的多项人工指标；近期价值更像沟通与知识辅助。它不支持“通过考试即获得临床资格”，也不支持“长答案更详细就更安全”。
\end{importantbox}

\subsection{本章小结}

Med-PaLM 在科学共识、推理、帮助程度等多项人类评价上显著优于 Flan-PaLM，并缩小与医生答案的差距；同时，详细回答可能引入更多错误，伤害与偏见仍需专家判断。最可信的课堂落点是受监督的临床增强，而不是自动替代。

